% Experiment report. Fixed sections: see experiments/TEMPLATE/report.tex and experiments/README.md.
% Build: npm run experiment:report -- 02-research-usefulness
\documentclass{../polyexp}
\begin{document}
% generated by scripts/experiment-report.ts from experiment.json and models/<run>/ - do not edit
\providecommand{\ExpHeader}{}\renewcommand{\ExpHeader}{%
{\sffamily\footnotesize\color{expgrey}EXPERIMENT 2 \textperiodcentered\ MECHANISM \hfill 2026-11-24}\par
{\LARGE\bfseries\sffamily Score research by usefulness}\hfill\verdictbox{expred}{NOT KEPT}\par\smallskip
{\itshape Does the bot get stronger when it prefers techs it can use right away?}\par\medskip
\begin{tabular}{@{}p{22mm}p{\dimexpr\linewidth-26mm}@{}}
\fact{Change}{The research score gains techUse per action the tech would unlock right now, and loses techUseless when it unlocks nothing.}
\fact{Knobs}{\knob{techUse} $\in$ \{3\}, best 3; \knob{techUseless} $\in$ \{15\}, best 12.8}
\fact{Base}{\model{greedy-full-v1}}
\fact{Method}{Tuning run (hill climbing), 100,000 games, against a control run of the same budget; metric: points share}
\fact{Setups}{4 players, 16$\times$16, drylands, turn cap 150}
\fact{Record}{issue \issue{29}, pull request \issue{36}; status: done}
\end{tabular}\par\medskip
\begin{center}\begin{tikzpicture}[node distance=5mm]
    \node[flowbox] (n0) {{\sffamily\bfseries\scriptsize\color{expgrey}BASE}\\[1pt]\model{full-v1}};
    \node[flowbox,right=5mm of n0] (n1) {{\sffamily\bfseries\scriptsize\color{expgrey}CHANGE}\\[1pt]\knob{techUse}\\{\scriptsize 3}\\\knob{techUseless}\\{\scriptsize 15}};
    \node[flowbox,right=5mm of n1] (n2) {{\sffamily\bfseries\scriptsize\color{expgrey}MEASURED}\\[1pt]Tuning run, 100,000 games, vs a control\\{\scriptsize 1 setup, metric: points share}};
    \node[flowbox,right=5mm of n2] (n3) {{\sffamily\bfseries\scriptsize\color{expgrey}RESULT}\\[1pt]Points vs base, final test\\{\scriptsize 50.0 $\rightarrow$ 39.8\,\%}\\\textbf{$-$10.2\pp}};
    \node[flowbox,right=5mm of n3] (n4) {{\sffamily\bfseries\scriptsize\color{expgrey}VERDICT}\\[1pt]\textbf{\color{expred}NOT KEPT}};
    \draw[flowarrow] (n0) -- (n1); \draw[flowarrow] (n1) -- (n2); \draw[flowarrow] (n2) -- (n3); \draw[flowarrow] (n3) -- (n4);
\end{tikzpicture}\end{center}
\begin{center}\sffamily\small\begin{tabular}{@{}lrrr@{}}\toprule
Measure & Base & New & Change\\\midrule
Points vs base, final test & 50.0\,\% & 39.8\,\% & $-$10.2\\
Won vs base, final test & 50.0\,\% & 28.2\,\% & $-$21.8\\
Points share, 6-model panel & 67.7\,\% & 63.3\,\% & $-$4.4\\
\bottomrule\end{tabular}\end{center}
}
\providecommand{\ExpResults}{}\renewcommand{\ExpResults}{%
\subsection{Tuning run, 100,000 games {\normalfont\small(\model{greedy-full-research})}}
\begin{center}\begin{tikzpicture}\begin{axis}[expplot,height=4.6cm,xlabel={training games (thousands)},ylabel={Points share, panel mean (\%)},enlarge y limits=0.2]
\addplot[dashed,expgrey,line width=0.7pt,forget plot] coordinates {(0.0,63.40) (100.0,63.40)};
\addplot[expblue,mark=*,mark size=1.8pt] coordinates {(0.0,63.40) (12.9,63.30) (25.2,63.30) (37.7,63.30) (50.3,63.30) (62.9,63.30) (75.2,63.30) (87.7,63.30) (100.0,63.30)};
\end{axis}\end{tikzpicture}\\[-2pt]{\footnotesize\color{expgrey}Learning curve of the tuning run (1 kept changes in 244 rounds). Dashed: the start, i.e. the mechanism at its hand-set value.}\end{center}
\begin{center}\sffamily\small\begin{tabular}{@{}lrr@{}}\toprule
Opponent & Won (\%) & Points (\%)\\\midrule
\model{v1} & 84 $\pm$ 4 & 69.6\\
\model{hill-v1} & 79 $\pm$ 4 & 65.0\\
\model{hill-v2} & 81 $\pm$ 4 & 65.0\\
\model{hill-v3} & 76 $\pm$ 4 & 60.8\\
\model{hill-v4} & 76 $\pm$ 4 & 60.8\\
\model{full-v1} & 28 $\pm$ 4 & 39.8\\
\midrule
\textbf{Mean} & 70.8 & 60.2\\
\bottomrule\end{tabular}\\[3pt]{\footnotesize\color{expgrey}Final test of the saved model: 400 unseen games against each panel model.}\end{center}
}

\ExpHeader

\section{Question and hypothesis}
Research is scored as a fixed weight minus cost and tier. It never asks whether the tech unlocks anything the player can use, so bots buy Fishing with no water in reach. \textbf{Hypothesis:} scoring a tech by what it would let the player do now avoids wasted stars and raises strength.

\section{Mechanism}
$\mathit{uses}(t)$ counts the actions tech $t$ would add right now: builds and harvests on own tiles it unlocks, unit types a city could train, abilities with a valid target. The research score gets $+\,\knob{techUse}\cdot\mathit{uses}(t) - \knob{techUseless}\cdot[\mathit{uses}(t)=0]$. Both knobs default to 0.

\section{Setup}
Same protocol as experiment 1: one hill-climbing run of 100,000 games from the base against the same control. Hand-set start: \knob{techUse} 3, \knob{techUseless} 15. The research score these compete with is about 28 minus cost and tier, so the start values were large against it.

\section{Results}
\ExpResults

The mechanism costs 4.3 points on the panel at its hand-set values (63.4\,\% against 67.7\,\%) and loses clearly to the base: 39.8\,\% of points, 28\,\% won ($\pm$4). Training could not repair it: one kept change in 244 rounds (\knob{techUseless} 15 to 12.8).

\section{Behaviour}
The behaviour metric planned for this step was not measured; the strength result was judged clear enough on its own. This is a gap in the record. Planned metric: share of researched techs whose unlocks are never used.

\section{Interpretation}
``What can I build right now'' undervalues techs whose payoff comes later (a unit tech before a war, the parent of a needed tech) and overvalues cheap terrain builds. The base had already tuned its research order through per-tech weights, so the new term mostly overrode a tuned order. The start was also about five times too strong, and multiplicative nudges of $\pm$40\,\% cannot walk a knob that far within the budget. This experiment is the reason later sweeps always include a weak value.

\todo{jakob}{Your reading: does this match what you saw? Anything to add or correct?}

\section{Threats to validity}
\begin{itemize}
\item Only one strength was effectively tested. A mild version (for example \knob{techUse} 0.3, \knob{techUseless} 2) is untested and could behave differently.
\item The behaviour metric was not measured.
\end{itemize}

\section{Verdict and next step}
\textbf{Not kept.} The code stays, switched off. A retry at a much weaker strength would be a new experiment. Next: front pull (experiment 3).

\section{Notes for the paper}
Draft: a myopic usefulness heuristic for research loses 10 points of share against tuned per-tech weights; and hill climbing cannot recover from a badly sized start. Motivates the move from tuning runs to strength sweeps.

\todo{jakob}{Is this worth a place in the paper? Where?}

\end{document}
